\slajdid{09_1-s025}
\begin{frame}[label=09_1-s025]
\gusttekst
Međutim očekujući da je ova tačka na polovini napona napajanja, „tako bi trebalo da se projektuju logička kola“, možemo u članu koji „nam smeta“ da to pretpostavimo. Na primer
\[\frac{V_{DD}-\frac{V_{DD}}2}{R_L}=\frac{k_n}{2}\frac{L_nE_{Cn}}{L_nE_{Cn}+\left(\frac{V_{DD}}2-V_{Tn}\right)}(V_S-V_{Tn})^2\]
\[V_S-V_{Tn}=\sqrt{\frac{V_{DD}\left(L_nE_{Cn}+\left(\frac{V_{DD}}2-V_{Tn}\right)\right)}{k_nR_LL_nE_{Cn}}}=\alert{0.26V\ \Rightarrow\ V_S=0.71V}\]
Za standardne vrednosti a tačan rezultat je \alert{0.7V}.\par\medskip
Greška koju smo napravili u odnosu na tačan rezultat je mala.\par\medskip
Razlog je što je stvarno tačka $V_S=0.7V$ dosta bliska polovini napona napajanja $\frac{V_{DD}}2=0.6V$.\par\bigskip
\alert{U svakom slučaju na ispitu je drugi, ovaj, način dozvoljen i poželjan. Prvi treba zaobilaziti u širokom luku pošto je jako lako napraviti greške u izvođenju izraza i na ispitu vrlo verovatne.}
\end{frame}
